% =====================================================================
%  config/settings.tex — gói lệnh và định dạng chung
% =====================================================================
\usepackage{iftex}

% ---------------- Phông chữ và tiếng Việt ----------------------------
\ifPDFTeX
  % pdfLaTeX: cần gói vntex (có sẵn trong TeX Live đầy đủ / Overleaf)
  \usepackage[T5]{fontenc}
  \usepackage[utf8]{inputenc}
  \usepackage{times}
\else
  % XeLaTeX / LuaLaTeX: TeX Gyre Termes (bản Times đầy đủ dấu tiếng Việt)
  \usepackage{fontspec}
  \setmainfont{texgyretermes}[
    Extension      = .otf,
    UprightFont    = *-regular,
    BoldFont       = *-bold,
    ItalicFont     = *-italic,
    BoldItalicFont = *-bolditalic]
  % Phông chữ đơn cách cho mã nguồn: TeX Gyre Cursor (kiểu Courier,
  % đủ dấu tiếng Việt)
  \setmonofont{texgyrecursor}[
    Extension      = .otf,
    UprightFont    = *-regular,
    BoldFont       = *-bold,
    ItalicFont     = *-italic,
    BoldItalicFont = *-bolditalic]
\fi
\IfFileExists{vietnamese.ldf}
  {\usepackage[vietnamese]{babel}}
  {\usepackage{babel}\babelprovide[import,main]{vietnamese}}

% ---------------- Trang ----------------------------------------------
\usepackage[a4paper,left=2.5cm,right=1.5cm,top=2.5cm,bottom=2.5cm,
            footskip=1.2cm]{geometry}
\usepackage{setspace}
\setstretch{1.3}
\setlength{\parindent}{1cm}
\setlength{\parskip}{4pt}
\usepackage{indentfirst}      % thụt lề cả đoạn đầu sau tiêu đề (kiểu VN)
\pagestyle{plain}

% ---------------- Toán học -------------------------------------------
\usepackage{amsmath,amssymb,amsthm}
\usepackage{mathtools}
\allowdisplaybreaks[2]

% ---------------- Bảng, hình, màu ------------------------------------
\usepackage{graphicx}
\graphicspath{{images/}}
\usepackage{array,booktabs,multirow,longtable}
\usepackage{float}
\usepackage[table]{xcolor}
\usepackage{tikz}
\usetikzlibrary{arrows.meta,positioning,shapes.geometric,calc}
\usepackage[font=normalsize,labelfont=bf,labelsep=period,
            justification=centering,skip=6pt]{caption}
\captionsetup[table]{position=top}
\numberwithin{equation}{chapter}
\numberwithin{figure}{chapter}
\numberwithin{table}{chapter}
\renewcommand{\arraystretch}{1.25}

% ======================================================================
%  QUY TẮC TRÌNH BÀY CHUNG (áp dụng cho mọi chương)
%  1. Cấp tiêu đề: CHƯƠNG > x.y > \mucnho{1) ...} > \mucphu{...} > nhãn
%     nghiêng đầu mục ("Khái niệm:") ngay trong dòng.
%     \mucnho đánh số 1), 2) khi mục có từ hai nội dung song song; không
%     đánh số khi chỉ là một khối riêng. Mọi tiêu đề đi liền đoạn sau.
%  2. Danh sách: cấp 1 "•", cấp 2 "–"; công thức/bảng thuộc ý nào thì
%     nằm TRONG mục danh sách của ý đó, không để rời giữa hai danh sách.
%  3. Định lý/Mệnh đề/Hệ quả: đầu đề đậm, thân nghiêng. Định nghĩa/Ví dụ:
%     đầu đề đậm, thân đứng; Ví dụ kết thúc bằng dấu ◊. Nhận xét/Lưu ý:
%     đầu đề nghiêng, thân đứng. Kết luận mở đầu bằng "→" = \begin{ketluan}.
%     Không tô nền, không vạch màu; chỉ "Quy ước chung" có khung đen mảnh.
%  4. Công thức: chỉ đánh số công thức cốt lõi hoặc được dẫn chiếu; phép
%     tính trung gian không đánh số. Mọi công thức đứng cùng câu dẫn.
%  5. Bảng: tiêu đề phía trên, đánh số "Bảng n.m."; kẻ ô đen mảnh; hàng
%     tiêu đề chữ đậm; chữ \small. Bảng nhỏ minh họa trong ví dụ thì
%     không đánh số.
%  6. Không để tiêu đề cuối trang, không tách công thức khỏi câu dẫn.
% ======================================================================
\usepackage{needspace}
% Tiêu đề cấp 3: "1) ...", riêng một dòng, đậm
\newcommand{\mucnho}[1]{\par\addvspace{8pt}\Needspace{5\baselineskip}%
  \noindent\textbf{\boldmath #1}\par\nopagebreak\vspace{2pt}}
% Tiêu đề cấp 4: riêng một dòng, đậm nghiêng
\newcommand{\mucphu}[1]{\par\addvspace{5pt}\Needspace{4\baselineskip}%
  \noindent\textbf{\boldmath\textit{#1}}\par\nopagebreak\vspace{1pt}}

% ---------------- Danh sách ------------------------------------------
\usepackage{enumitem}
\setlist{itemsep=3pt,topsep=3pt,parsep=0pt,partopsep=0pt}
\setlist[itemize,1]{label=\textbullet,leftmargin=1.0cm,labelsep=0.5em}
\setlist[itemize,2]{label=--,leftmargin=0.9cm,labelsep=0.5em}
\setlist[itemize,3]{label=$\circ$,leftmargin=0.9cm,labelsep=0.5em}
\setlist[enumerate]{leftmargin=1.0cm,labelsep=0.5em}

% ---------------- Mã nguồn -------------------------------------------
% Dùng minted (tô màu cú pháp, cần biên dịch với -shell-escape và Python
% Pygments; Overleaf có sẵn). Nếu không bật shell-escape, tự chuyển sang
% fvextra (không tô màu). Gói listings KHÔNG dùng vì in sai thứ tự chữ
% tiếng Việt (ồ, ữ, ơ, đ...) khi biên dịch bằng XeLaTeX.
\usepackage{fvextra}
\newif\ifdungminted
\ifdefined\shellescape\ifnum\shellescape=1 \dungmintedtrue\fi\fi
\ifdefined\pdfshellescape\ifnum\pdfshellescape=1 \dungmintedtrue\fi\fi
\ifdungminted
  \usepackage{minted}
  \setminted{style=friendly,fontsize=\footnotesize,linenos,numbersep=6pt,
    frame=single,framesep=2mm,rulecolor=\color{gray!60},
    breaklines=true,tabsize=4}
\fi
\fvset{fontsize=\footnotesize,frame=single,framesep=2mm,
  rulecolor=\color{gray!60},breaklines=true,tabsize=4}
% Kiểu chú thích "Mã x.y" và danh sách mã nguồn
\DeclareCaptionType[fileext=lol,within=chapter]{chuongtrinh}[Mã][Danh sách mã nguồn]
% \maNguon{đường dẫn file}{chú thích}{nhãn}
\newcommand{\maNguon}[3]{%
  \par\medskip\noindent
  \begin{minipage}{\linewidth}\captionof{chuongtrinh}{#2}\label{#3}\end{minipage}
  \par\nopagebreak
  \ifdungminted
    \inputminted{python}{#1}%
  \else
    \VerbatimInput[numbers=left,numbersep=6pt]{#1}%
  \fi
  \medskip}

% ---------------- Khung nhắc nội dung (ẩn khi nộp) -------------------
\usepackage[most]{tcolorbox}
\newif\ifhienkhung
\newcommand{\canviet}[1]{%
  \ifhienkhung
  \begin{tcolorbox}[breakable,enhanced,colback=gray!6,colframe=gray!55,
      boxrule=0.4pt,arc=1pt,left=6pt,right=6pt,top=3pt,bottom=3pt,
      fonttitle=\small\bfseries,coltitle=black,colbacktitle=gray!18,
      title={Nội dung cần viết}]
    \small\itshape #1
  \end{tcolorbox}
  \fi}
% Khung "quy ước" — là nội dung thật của báo cáo, luôn hiện:
% khung đen mảnh, tiêu đề đậm ngay trong khung
\newtcolorbox{quyuoc}[1][Quy ước]{breakable,colback=white,colframe=black,
  boxrule=0.5pt,arc=0pt,left=7pt,right=7pt,top=5pt,bottom=5pt,
  before skip=8pt,after skip=10pt,
  before upper={\textbf{#1}\par\smallskip}}

% ---------------- Định lý, định nghĩa, ví dụ -------------------------
\usepackage{etoolbox}
\theoremstyle{plain}
\newtheorem{dinhly}{Định lý}[chapter]
\newtheorem{menhde}[dinhly]{Mệnh đề}
\newtheorem{hequa}[dinhly]{Hệ quả}
\theoremstyle{definition}
\newtheorem{dinhnghia}{Định nghĩa}[chapter]
\newtheorem{vidu}{Ví dụ}[chapter]
\theoremstyle{remark}
\newtheorem{nhanxet}{Nhận xét}[chapter]
\newtheorem*{luuy}{Lưu ý}
\renewcommand{\proofname}{Chứng minh}
% Dấu ◊ kết thúc mỗi ví dụ
\AtEndEnvironment{vidu}{\unskip\nobreak\hfill\ensuremath{\diamond}}
% Kết luận mở đầu bằng mũi tên (không khung, không nền)
\newenvironment{ketluan}
  {\par\addvspace{4pt}\begingroup\leftskip=1.4em\noindent
   \hspace*{-1.4em}\makebox[1.4em][l]{$\rightarrow$}\ignorespaces}
  {\par\endgroup\addvspace{4pt}}

% ---------------- Tên các mục (tiếng Việt) ---------------------------
\addto\captionsvietnamese{%
  \renewcommand{\contentsname}{MỤC LỤC}%
  \renewcommand{\listfigurename}{Danh sách hình vẽ}%
  \renewcommand{\listtablename}{Danh sách bảng}%
  \renewcommand{\figurename}{Hình}%
  \renewcommand{\tablename}{Bảng}%
  \renewcommand{\bibname}{TÀI LIỆU THAM KHẢO}%
  \renewcommand{\chaptername}{CHƯƠNG}%
  \renewcommand{\appendixname}{Phụ lục}%
}

% ---------------- Tiêu đề chương, mục --------------------------------
\usepackage{titlesec}
\newcommand{\chuongfont}{\fontsize{13pt}{17pt}\selectfont\bfseries}
\titleformat{\chapter}[block]
  {\chuongfont\boldmath\centering}{CHƯƠNG \thechapter.}{1em}{}
\titleformat{name=\chapter,numberless}[block]
  {\chuongfont\centering}{}{0pt}{\MakeUppercase}
\titlespacing*{\chapter}{0pt}{0pt}{18pt}
\titleformat{\section}
  {\normalsize\bfseries\boldmath}{\thesection.}{0.8em}{}
\titlespacing*{\section}{0pt}{12pt}{6pt}
\titleformat{\subsection}
  {\normalsize\bfseries\boldmath}{\thesubsection.}{0.8em}{}
\titlespacing*{\subsection}{0pt}{10pt}{4pt}
\titleformat{\subsubsection}
  {\normalsize\bfseries\itshape}{\thesubsubsection.}{0.8em}{}
\titlespacing*{\subsubsection}{0pt}{8pt}{3pt}
% Không để tiêu đề mục (x.y, x.y.z) nằm lẻ loi cuối trang
\usepackage{etoolbox}
\pretocmd{\section}{\Needspace{6\baselineskip}}{}{}
\pretocmd{\subsection}{\Needspace{5\baselineskip}}{}{}
\setcounter{secnumdepth}{3}
\setcounter{tocdepth}{1}       % mục lục đến cấp x.y như khung

% ---------------- Mục lục, danh sách hình, bảng ----------------------
\usepackage[titles]{tocloft}
\renewcommand{\cftdotsep}{2}
% Cấp chương
\renewcommand{\cftchapfont}{\bfseries}
\renewcommand{\cftchappagefont}{\bfseries}
\renewcommand{\cftchappresnum}{CHƯƠNG~}
\renewcommand{\cftchapaftersnum}{.}
\setlength{\cftchapnumwidth}{2.75cm}
\renewcommand{\cftchapleader}{\bfseries\cftdotfill{\cftdotsep}}
\setlength{\cftbeforechapskip}{9pt}
% Cấp mục
\setlength{\cftsecindent}{1.0cm}
\setlength{\cftsecnumwidth}{1.1cm}
\renewcommand{\cftsecaftersnum}{.}
\setlength{\cftbeforesecskip}{1pt}
% Hình, bảng
\setlength{\cftfigindent}{0pt}
\setlength{\cftfignumwidth}{1.5cm}
\renewcommand{\cftfigpresnum}{Hình~}
\renewcommand{\cftfigaftersnum}{.}
\setlength{\cftfignumwidth}{2.1cm}
\setlength{\cfttabindent}{0pt}
\renewcommand{\cfttabpresnum}{Bảng~}
\renewcommand{\cfttabaftersnum}{.}
\setlength{\cfttabnumwidth}{2.1cm}
% Tiêu đề trang mục lục / danh sách: in hoa, giữa trang
\makeatletter
\renewcommand{\tableofcontents}{%
  \chapter*{MỤC LỤC}\@starttoc{toc}}
\renewcommand{\listoffigures}{%
  \chapter*{DANH SÁCH HÌNH VẼ}\@starttoc{lof}}
\renewcommand{\listoftables}{%
  \chapter*{DANH SÁCH BẢNG}\@starttoc{lot}}
\makeatother

% Chương không đánh số nhưng có trong mục lục (Tóm tắt, Lời cảm ơn, ...)
\newcommand{\chuongphu}[1]{%
  \cleardoublepage\phantomsection
  \addcontentsline{toc}{chapter}{#1}%
  \chapter*{#1}}

% ---------------- Liên kết -------------------------------------------
\usepackage[unicode,hidelinks,bookmarksnumbered]{hyperref}
\hypersetup{
  pdftitle={Mã Hamming (7,4): mô phỏng mã hóa, gây nhiễu, phát hiện và sửa lỗi cho chuỗi ký tự},
  pdfauthor={Nhóm 10 – Lớp L06 – MT1007},
  pdfsubject={Báo cáo bài tập lớn Đại số tuyến tính}}
